\begin{abstract}
We compare two preference-optimization methods, Direct Preference Optimization (DPO) and Kahneman--Tversky Optimization (KTO), as a way to align a language model to write humor about current events.  We fine-tune Mistral-7B-Instruct-v0.2 with from-scratch implementations of both losses on comedy prompts built from Guardian news articles, and vary the prompt context length: a short condition (today's headline and lede only) and a long condition that also includes a model-generated backstory summarizing recent related coverage.  This gives four trained variants.  Evaluated against the base model by an LLM cross-judge on a held-out set of 30 prompts, short-context training wins 62\% of comparisons, while long-context training wins only 47\% (worse than no fine-tuning at all).  The long-context KTO variants reach the lowest validation losses of the four runs but the cross-judge ranks them no better than base, and instrumenting the loss shows why: their mean policy-vs-reference log-probability ratio collapses to large negative values during training, indicating the policy has drifted far from the reference, while KTO's bounded loss never reflects that drift.  The practical message: under bounded surrogate losses, validation loss alone is not a reliable training signal, and longer contexts can hurt rather than help.
\end{abstract}
